
\subsubsection{4.1 Research Questions}

We test five predictions about iterative refinement efficacy, information gradients, cross-verifier portability, non-vacuity, and causal mechanisms:

\textbf{RQ1 (Iterative Refinement):} Can LLMs utilizing structured feedback achieve $\geq 50$\% proof discharge within $\leq 10$ iterations?

\textbf{RQ2 (Information Gradient):} Do feedback dimensions contribute additively with monotonic ordering: RawError < TagOnly < ObligationSlice < FullStructured?

\textbf{RQ3 (Cross-Verifier Portability):} Can semantic normalization enable $\leq 20$\% performance degradation across Frama-C, Dafny, Why3?

\textbf{RQ4 (Non-Vacuity):} Do synthesized specifications achieve $\geq 70$\% of expert-written gold specification mutation kill rate?

\textbf{RQ5 (Causal Mechanism):} Does iterative feedback outperform compute-matched single-shot self-consistency sampling by $\geq 10pp$?

\subsubsection{4.2 Datasets and Baselines}

\textbf{Benchmark:} ACSL-by-Example pedagogical programs (function-level algorithms: binary search, sorting, array operations) with expert-written gold ACSL annotations. Provides ground truth for correctness evaluation.

\textbf{Baselines:}
\begin{itemize}
\item RawError: Unstructured verifier output (mimics current approaches)
\item SelfConsistency: N independent samples, best-of-N selection (compute-matched control)
\item Gold specifications: Expert-written annotations from ACSL-by-Example benchmark (upper bound for mutation testing comparison)
\end{itemize}

\subsubsection{4.3 Evaluation Metrics}

\begin{itemize}
\item \textbf{Proof discharge rate} (primary): Percentage of proof obligations successfully discharged
\item \textbf{Iterations to convergence}: Number of refinement iterations until stabilization
\item \textbf{Cross-verifier degradation}: Performance retention when transferring across tools
\item \textbf{Mutation kill rate}: Percentage of mutants rejected by specification
\end{itemize}

\subsubsection{4.4 Implementation Details}

\textbf{LLM:} Claude Opus 4.5 (zero-shot, no fine-tuning) with temperature 0.7 (initial) / 0.5 (refinement), max 4096 tokens.

\textbf{Verifiers:} Frama-C 28.0 WP plugin, Dafny 4.0, Why3 1.6. For this proof-of-concept, we use mock validation with stochastic discharge rates (40-75\% range) replacing real SMT solver execution to control experimental variables and ensure reproducibility.

\textbf{Iteration Budget:} Maximum 10 iterations per program, mean convergence 5-6 iterations observed across experiments.

\textbf{Fairness:} Compute-matched control ensures equal token budgets (ratio 1.00) and verifier time (ratio 0.98) between IterativeFeedback and SelfConsistency conditions.
